\section{Architecture and length boundaries}
\label{sec:5}
\subsection{Fixed Protenix adapters beyond their adaptation lengths}\label{fixed-protenix-adapters-beyond-their-adaptation-lengths}

On Length48, native full-factor adapters trained on Train384 score 0.52919,
versus 0.48667 for trained rotations and 0.32122 for query-only. The locked
primary difference is \textbf{+0.04252 {[}0.02899, 0.05765{]}}, with 41/48 positive
target differences. Residue-averaged lDDT gives +0.04029 and TM-score +0.04782,
both with positive intervals. Native-query (+0.20797) is a descriptive secondary
comparison. All 672 model--target predictions complete without input cropping.

Shared residue updates therefore remain useful beyond the adapter's training
lengths in this test. The experiment does not establish arbitrary-length,
family or pretraining generalization. The three length-bin effects are positive,
but do not establish a monotonic length trend. Protenix G+ was not evaluated
on this longer-chain panel; its AF2 evaluation below is a different experiment.

\subsection{AF2: supported effects and competitive boundaries}\label{af2-supported-effects-and-competitive-boundaries}

\begin{table}[t]
\centering
\caption{AF2/OpenFold at the fixed 1,536-update endpoint. Both panels were previously observed. Means average three training seeds and, for rotated models, three rotations. G+ remains in the main comparison regardless of ranking.}
\label{tab:openfold}
\small\setlength{\tabcolsep}{3pt}
\begin{tabular*}{\linewidth}{@{\extracolsep{\fill}}llrrrr@{}}
\toprule
Adaptation & Panel & Query & Native & Rotated & G+ \\
\midrule
Train96 & Confirm96-B & .30191 & .50190 & .48209 & .49534 \\
Train96 & Length48 & .26079 & .38872 & .37834 & .39849 \\
Train384 & Confirm96-B & .30191 & .50684 & .50381 & .51333 \\
Train384 & Length48 & .26079 & .40165 & .39040 & .40868 \\
\midrule
\multicolumn{2}{l}{Native minus rotation} & \multicolumn{4}{r}{Mean [95\% target-bootstrap interval]} \\
Train96 & Confirm96-B & \multicolumn{4}{r}{$+.01981\ [+.00928,+.03129]$} \\
Train96 & Length48 & \multicolumn{4}{r}{$+.01038\ [+.00403,+.01705]$} \\
Train384 & Confirm96-B & \multicolumn{4}{r}{$+.00303\ [-.00675,+.01218]$} \\
Train384 & Length48 & \multicolumn{4}{r}{$+.01126\ [+.00527,+.01734]$} \\
\bottomrule
\end{tabular*}
\end{table}


All three adapter families improve mean scores over query-only. Direction
effects are supported on both Train96 panels and on Length48 after Train384
adaptation. At Train384 on Confirm96-B, all three metric intervals cross zero;
one of three training seeds favors rotation. This establishes neither a native
advantage nor equivalence in that condition. The longer-chain Train384 effect
has positive intervals for both supplementary metrics, although one TM-score
rotation marginal is slightly negative.

A post-result paired analysis finds Train384 Native-G+ of
-0.00648 {[}-0.01777, 0.00536{]} on Confirm96-B and
-0.00703 {[}-0.01369, -0.00039{]} on Length48. Train96 Native-G+ on Length48
is also negative, -0.00977 {[}-0.01674, -0.00288{]}. These unadjusted intervals
support a longer-chain competitive limitation in the tested recipes. They do
not show that G+ learned a particular alignment mechanism.

The direction-effect change from Train96 to Train384 is
-0.01678 {[}-0.03214, -0.00261{]} on Confirm96-B and
+0.00087 {[}-0.00728, 0.00863{]} on Length48. This directly contrasts the paired
effects; it does not infer an interaction from separate significance decisions.
The short-panel pattern differs from Protenix. Neither interaction identifies
which architectural or optimization property causes that difference.

\subsection{AtlasFold: fixed non-OPM extension slot}\label{atlasfold-fixed-non-opm-extension-slot}

\textbf{Pending complete matrix; no adapter-effect claim.} This extension asks whether
native orientation matters at a product/difference interface while retaining
AtlasLM-3B. The locked matrix is Train96, 1,536 updates, three native fits,
three rotations crossed with three seeds, and three G+ fits. Symmetric Dev8
calibration has selected the common \(10^{-4}\) learning rate. The primary
comparison is native-rotation on observed Confirm96-B; Length48 and G+ are
secondary. Table 1 already specifies the construction independently of results.

The final table will report query-only, native, rotation and G+ scores on both
panels, paired intervals, all seed marginals and prediction failures. A positive,
null or negative outcome fills the same slot and does not change the main
question. Completed original AtlasFold predictions in Table 4 do not count as
this adapter experiment.
